\section{Fixed-branch irreducibility for every vertical fiber}
\label{sec:irreducibility}

This section closes the algebraic premise needed in
Section~\ref{sec:schinzel}.  Retain
\[
 A=1+4a^2,\qquad s=\frac{a-1}{2},\qquad
 Z=z^3,\qquad D=1-Z-a^2Z^2,
\]
\[
 N_g(b)=16a^4b^2-A(b-1)^4
\]
and
\begin{equation}
 P(b)=16D^2A^2b^4-\delta_\tau a^4Z^4N_g(b)^2
       -32A^3s^2D^2b^5. \label{eq:value-P}
\end{equation}
% src: data/l22_elimination.jsonl (integral-model row)
The class construction uses only the fixed canonical branch
\[
 \tau^\dagger=\frac{1+2a^2}{A},
 \qquad
 \delta_{\tau^\dagger}=-\frac{4a^4}{A}.
\]
On this branch multiplication by the nonzero scalar $A/4$ gives
\begin{equation}
 H:=\frac A4P
 =a^8Z^4N_g(b)^2+4A^3D^2b^4
  -2A^4(a-1)^2D^2b^5.                 \label{eq:integral-model}
\end{equation}
Its constant and leading coefficients are equal:
\begin{equation}
 H_0=H_8=L=a^8A^2Z^4.                 \label{eq:equal-endpoints}
\end{equation}
% src: data/l22_elimination.jsonl (integral-model row)

\subsection{The off-branch warning}

The branch qualification is essential.  If $a=1$ and
$\delta_\tau=\sigma^2\in\Q^{\times2}$, then
\[
 P(b)=
 \bigl(4DAb^2-\sigma a^2Z^2N_g(b)\bigr)
 \bigl(4DAb^2+\sigma a^2Z^2N_g(b)\bigr),
\]
a genuine quartic-by-quartic factorization whenever the displayed
endpoints are nonzero.  This does not meet the canonical branch, because
$\delta_{\tau^\dagger}=-4a^4/A<0$ for rational $a\ne0$.
% src: data/l21_reducible_locus.jsonl (l21a and l21b rows)
Thus no blanket assertion over a moving $\tau$ is used below: every
irreducibility statement is about the fixed pair
$(\tau^\dagger,\delta_{\tau^\dagger})$.

\subsection{Endpoint and Newton elimination}

\begin{theorem}[Fixed-\texorpdfstring{$Z$}{Z} canonical-branch
irreducibility; PROVED]
\label{thm:fiber-irred}
Fix $Z\in\Q^\times$.  On
\[
 \tau=\tau^\dagger=\frac{1+2a^2}{1+4a^2},
 \qquad
 \delta=-\frac{4a^4}{1+4a^2},
\]
the polynomial $H$ of \eqref{eq:integral-model}, and hence $P$, is
irreducible of degree eight in $\Q(a)[b]$.
% src: data/l22_elimination.jsonl (fixed-Z-theorem row)
\end{theorem}

\begin{proof}
We first assume $Z\notin\{0,1\}$.  The polynomial $H$ is primitive in
$\Q[a][b]$.  Indeed, a common divisor of its coefficients must divide the
equal endpoints $a^8A^2Z^4$.  Since $A=1+4a^2$ is irreducible in $\Q[a]$
and coprime to $a$, only $a$ and $A$ could occur.  The coefficient of
$b^4$ at $a=0$ is $4(1-Z)^2\ne0$, excluding $a$, and direct reduction of
that coefficient modulo $A$ excludes $A$.  Gauss's lemma therefore turns
any factorization over $\Q(a)$ into one by primitive factors in
$\Q[a,b]$.
% src: data/l22_elimination.jsonl (integral-model and local-newton-data rows)

At the prime $A$, the coefficient valuations in ascending powers of $b$
are
\[
 (2,2,1,1,0,1,1,2,2).
\]
The Newton polygon consequently has two sides, each of horizontal length
four, with slopes $-1/2$ and $+1/2$.  The slopes belonging to a global
factor form a sub-multiset of these slopes.  Its endpoint valuations are
integers, whereas the sum of an odd number of half-integral root
valuations is not.  Every global factor therefore has even degree.  If
the octic were reducible, it would have a factor of degree two or four.
% src: data/l22_elimination.jsonl (local-newton-data and global-degree-filter rows)

At $a=0$ the coefficient valuations are
\[
 (8,8,8,8,0,0,8,8,8),
\]
and the sides have data
\[
 4@(-2),\qquad 1@0,\qquad 3@(8/3).
\]
For the four roots in the first cluster, substitution $b=a^2x$ gives the
residual polynomial
\begin{equation}
 R_0(x)=Z^4+4(1-Z)^2x^4.               \label{eq:small-residual}
\end{equation}
It has no rational linear factor, since its two terms are nonnegative on
$\Q$ and do not vanish simultaneously.  It therefore has no rational
cubic factor either.  The denominator three in the last slope makes its
three roots an indivisible local block.  Hence a degree-two factor must
take two roots from the first four-root cluster; a degree-four factor
either takes all four of those roots or takes the unit root together with
the three-root cluster.  In a $4+4$ factorization one of the factors is
therefore the all-small factor.
% src: data/l22_elimination.jsonl (local-newton-data and global-degree-filter rows)

The two exact residual limits at infinity are
\begin{align}
 a^{-12}H(a,x/a)&\longrightarrow16Z^4(4x^2-1)^2, \label{eq:inf-small}\\
 a^{-20}H(a,ax)&\longrightarrow16Z^4x^4(x^2-4)^2. \label{eq:inf-large}
\end{align}
% src: data/l22_elimination.jsonl (local-newton-data row)
Let $F$ be a primitive factor of degree $r\in\{2,4\}$ formed from
$a=0$ small roots.  Unique factorization and
\eqref{eq:equal-endpoints} force
\[
 \operatorname{lc}_b(F)=cA^\beta,\qquad
 F(0)=d\,a^{2r}A^\delta,
\]
with $c,d\in\Q^\times$ and $0\le\beta,\delta\le2$.  Fixed powers of the
fixed rational $Z$ are units and are absorbed into $c,d$.  If $\ell$ of
the roots of $F$ are large at infinity, comparison with
\eqref{eq:inf-small}--\eqref{eq:inf-large} gives
\[
 2r+2\delta-2\beta=2\ell-r.
\]
The integral possibilities are
\[
\begin{array}{c|c}
r&(\beta,\delta,\ell)\\ \hline
2&(1,0,2),\ (2,0,1),\ (2,1,2),\\
4&(2,0,4).
\end{array}
\]
% src: data/l22_elimination.jsonl (global-degree-filter row)

Suppose first that $r=2$.  The rational degree-two divisors of the
repeated clusters in \eqref{eq:inf-small}--\eqref{eq:inf-large}, in the
three endpoint cases just listed, force the constant-to-leading ratio of
the $a=0$ residual factor to be
\[
 \rho=d/c\in\{16,-16\}.
\]
After monic normalization, divisibility of \eqref{eq:small-residual} has
the form
\[
 (x^2+ex+\rho)(x^2-ex+q)=x^4+K,\qquad
 K=\frac{Z^4}{4(1-Z)^2}>0.
\]
Coefficient comparison gives
$e(q-\rho)=0$ and $q-e^2+\rho=0$.  If $e=0$, then $q=-\rho$ and
$K=-\rho^2<0$.  Otherwise $q=\rho$ and $e^2=2\rho$, which asks for a
rational square root of $32$ or of $-32$.  Both alternatives are
impossible, so no quadratic factor exists.
% src: data/l22_elimination.jsonl (degree-2-elimination row)

Now let $r=4$.  The endpoint enumeration gives
$(\beta,\delta,\ell)=(2,0,4)$, so
\[
 \operatorname{lc}_b(F)=cA^2,\qquad F(0)=d\,a^8.
\]
The $b=ax$ initial form is the whole nonzero-root cluster
$16c(x^2-4)^2$, and hence $d=256c$.  At $b=a^2x$, the complementary
factor contributes the constant $Z^4/d$.  Matching the coefficient of
$x^4$ in \eqref{eq:small-residual} gives
\[
 \frac{cZ^4}{d}=4(1-Z)^2,\qquad
 Z^4=1024(1-Z)^2.                       \label{eq:forced-Z}
\]
The last equation splits into
\[
 Z^2+32Z-32=0,\qquad Z^2-32Z+32=0.
\]
Their discriminants are respectively
$1152=576\cdot2$ and $896=64\cdot14$, neither a rational square.
Thus no degree-four factor exists.
% src: data/l22_elimination.jsonl (degree-4-elimination row)

It remains to treat the degenerate residual value $Z=1$.  At $a=1$ the
primitive ascending coefficient vector of $H$ is
\[
 (25,-200,540,-760,1546,-760,540,-200,25).
\]
Modulo $11$, monic normalization gives
\[
 q(b)=1+3b+4b^2+7b^3+2b^4+7b^5+4b^6+3b^7+b^8.
\]
In $\F_{11}[b]/(q)$ the exact Rabin certificate is
\[
 b^{11^8}=b,\qquad
 b^{11^4}=8+7b+4b^2+9b^3+4b^4+7b^5+8b^6+10b^7,
\]
and $\gcd(b^{11^4}-b,q)=1$.  Since two is the only prime divisor of
eight, $q$ is irreducible.  Localizing $\Q[a]$ at
$L=a^8A^2$ makes $H/L$ monic; a putative factorization has monic factors
in that localization, and evaluation at $a=1$ preserves their positive
degrees because $L(1)=25\ne0$.  The certificate is therefore a
contradiction, proving the vertical statement at $Z=1$.
% src: data/l22_elimination.jsonl (exceptional-fiber Z=1 row)

Finally, $Z=0$ gives
\[
 H=2A^3b^4\bigl(2-A(a-1)^2b\bigr),
\]
which is reducible and is excluded because the cell parameter $z$ is
nonzero.  For fixed $Z\ne0$, the equation $D=0$ removes at most two
special values of $a$; it does not affect irreducibility in
$\Q(a)[b]$.  This completes all rational fixed fibers in the theorem.
% src: data/l22_elimination.jsonl (exceptional-fiber and fixed-Z-theorem rows)
\end{proof}

\subsection{Hilbert selection inside the admissible progression}

\begin{corollary}[Uniform admissible specialization; PROVED]
\label{cor:admissible-irred}
For every odd cell prime $w$ and every $z\in\Q^\times$ with
$v_w(z)\ge1$, the existing odd-$a$ character progression contains an
$a$ for which the canonical-branch polynomial $P(a,z^3;b)$ is
irreducible in $\Q[b]$ and $D\ne0$.
\end{corollary}

\begin{proof}
Fix $Z=z^3\ne0$.  Theorem~\ref{thm:fiber-irred} makes the octic over
$\Q(a)$ irreducible and, in characteristic zero, separable.  Quantitative
Hilbert irreducibility puts the integers at which its degree drops or it
becomes reducible in a thin set $E_Z$ satisfying
\begin{equation}
 \#\{n\in E_Z\cap\Z:|n|\le B\}
   =O_Z(\sqrt B\log B).                 \label{eq:thin-bound}
\end{equation}
\cite{cohenhit}\cite[Section~13, Theorems~1--2]{serremw}
% src: data/l22_elimination.jsonl (fixed-Z-theorem row);
% src: data/l21_irred_generic.jsonl (quantitative-Hilbert selection)

The character sum used in the class construction supplies
$r\bmod w$ with
\[
 \leg{1+4r^2}{w}=-1.
\]
Combining this residue with parity gives an odd
$a_0\bmod 2w$; every positive $a=a_0+2wk$ is character-admissible.
The progression contributes linearly many integers up to height $B$,
which eventually dominates \eqref{eq:thin-bound}.  It therefore contains
a point outside $E_Z$, and the equation $D=0$ excludes at most two more
points.  Such an $a$ has the asserted properties.
% src: data/l19_classexist.jsonl (character-sum rows);
% src: data/l21_irred_generic.jsonl (thin-set progression argument)
\end{proof}

This is the load-bearing order of the closure.  Fixed-$Z$ vertical
irreducibility gives Hilbert selection \emph{inside the already admissible
$a$ progression}; the resulting L20 class certificate has $f=w$ and all
Schinzel-pair conditions; the L19 fixed-pair implication applies classical
Schinzel~H to $\{q_1+Nt,G(t)\}$; the supplied prime pair gives an
emergent-free member; and the L6 equivalence yields the six-variable
definition.  No separate per-cell irreducibility premise remains.
% src: data/l22_elimination.jsonl (fixed-Z-theorem row)
% src: data/l20_admissible.jsonl; data/l18_schinzel_implies_h.jsonl

\subsection{Independent reciprocal-trace corroboration}

\begin{theorem}[The fixed reciprocal specialization; PROVED]
\label{thm:reciprocal-fixed}
On the same canonical branch, specialize
\[
 a=1,\quad A=5,\quad s=0,\quad
 \tau^\dagger=3/5,\quad\delta_{\tau^\dagger}=-4/5.
\]
For every $Z\in\Q^\times$, the resulting polynomial $P(1,Z;b)$ is
irreducible in $\Q[b]$.
% src: data/l22_reciprocal_cube.jsonl (meta and summary rows)
\end{theorem}

\begin{proof}
For this fixed choice
$(a,A,s,\tau,\delta)=(1,5,0,3/5,-4/5)$, put
$D=1-Z-Z^2$ and $u=b+b^{-1}$.  Direct expansion gives
\[
 P(b)=b^4T(u),\qquad
 T(u)=400D^2+\frac45Z^4\bigl(16-5(u-2)^2\bigr)^2.
\]
In $E=\Q(t)$, $t^2=-5$, one has
\[
 \frac54T=Q_-Q_+,\qquad
 Q_\pm=Z^2\bigl(16-5(u-2)^2\bigr)\pm10Dt.
\]
The conjugate quadratics are coprime, and trace reducibility is equivalent
to $80Z^2+50Dt$ being a square in $E$.  Its norm makes
\[
 125D^2+64Z^4
\]
a necessary rational square.  Writing $Z=p/q$ in lowest terms and
$d=q^2-pq-p^2$, the cleared numerator is
$125d^2+64p^4\equiv5\pmod8$: coprimality makes $d$ odd in all three
parity patterns.  It cannot be a square, so $T$ is irreducible.
% src: data/l22_reciprocal_cube.jsonl
%      (symbolic_identity and trace_reducibility rows)

For a root $u$ of $T$, the reciprocal lift is irreducible exactly when
$u^2-4$ is not a square in $\Q(u)$.  Its norm is a rational square times
\[
 \Xi(Z)=(125D^2+64Z^4)(125D^2+1024Z^4).
\]
Suppose $\Xi(Z)$ were a square and put $h=5D/(8Z^2)$.  The resulting point
on
\[
 v^2=(5h^2+1)(5h^2+16)
\]
maps by
\[
 x=25h^2,\qquad y=25hv
\]
to
\[
 E_0:\quad y^2=x(x+5)(x+80).
\]
% src: data/l22_reciprocal_cube.jsonl
%      (reciprocal_lift, Xi_theorem, and curve_maps rows)
The complete $2$-isogeny descent checks the square classes
\[
\begin{array}{c|c|c}
&\text{candidate classes}&\text{realized classes}\\ \hline
E_0&\{\pm1,\pm2,\pm5,\pm10\}&\{1,-5\}\\
E'_0:y^2=x(x-45)(x-125)&\{\pm1,\pm3,\pm5,\pm15\}&\{1,5\}.
\end{array}
\]
The excluded classes are eliminated modulo $3$, modulo $16$, or at the
real place, so the descent gives rank zero.  The good-prime point counts
$8$, $16$, and $12$ at $7$, $11$, and $13$ force torsion order at most
four; the visible points already give
\[
 E_0(\Q)=\{O,(0,0),(-5,0),(-80,0)\}.
\]
But $D\ne0$ for rational $Z$ and the image has $x=25h^2>0$, a
contradiction.  Thus $\Xi(Z)$ is nonsquare and the reciprocal lift is
irreducible.
% src: data/l22_reciprocal_cube.jsonl
%      (two_isogeny_descent, isogeny_cover, and elliptic_torsion rows)
\end{proof}

The value $a=1$ in Theorem~\ref{thm:reciprocal-fixed} is a
\emph{specialization witness} for the rational-function fiber.  It need
not satisfy $\leg{5}{w}=-1$ and is not the class parameter selected in
Corollary~\ref{cor:admissible-irred}.  Localizing at
$a^8A^2Z^4$ shows that this theorem also gives an independent proof of
fixed-$Z$ vertical irreducibility; the endpoint--Newton proof above remains
the primary closure.
% src: data/l22_reciprocal_cube.jsonl (branch_scope_note and summary rows)

\subsection{Route discipline}

Three structurally useful routes are deliberately not cited as the
closure.  The infinity calculation alone resolves four quadratic local
clusters but still permits global factor degrees $2$, $4$, and $6$; the
equal-endpoint comparison in the primary proof is what eliminates them.
A factor-tuple version of Schinzel controls only the product of the outside
Hilbert signs unless an all-plus compatibility class is also proved, so it
relocates rather than removes irreducibility.  Finally, the free-$\lambda$
square-branch theorem varies $\tau$ through the infinite L11c family; that
new branch parameter would cost a seventh witness and is excluded from the
six-count.  These are proved scope statements about alternate routes, not
gaps in Theorem~\ref{thm:fiber-irred}.
% src: data/l22_infinity.jsonl (verdict row)
% src: data/l22_factor_tuple.jsonl (summary row)
% src: data/l22_square_branch.jsonl (summary row)
